\chapter{Contexte du chip SPADMIC}
\label{chap:spadmic}
\justifying
\sloppy

Le projet \textbf{SPADMIC} s'inscrit dans une démarche d'intégration d'une électronique de lecture dédiée à une matrice de \textbf{SPAD} (\textit{Single-Photon Avalanche Diode}). Dans un tel système, l'information utile n'est pas limitée à la simple détection d'un événement : l'instant précis auquel cet événement se produit devient lui aussi une grandeur à exploiter. La mesure temporelle fine n'est donc pas un simple complément fonctionnel ; elle conditionne directement la valeur scientifique du système complet, la qualité de la discrimination des événements et la pertinence des traitements numériques réalisés en aval.

Dans le contexte de cette alternance, le bloc de conversion temps-numérique apparaît comme l'une des briques les plus structurées et les plus documentées de l'architecture. Le chip \textbf{SPADMIC} dans son ensemble reste en cours de consolidation, mais le \textbf{TDC} développé autour du projet \textbf{MPTDC} dispose déjà d'un socle théorique clair, d'une implémentation RTL complète, d'une stratégie de calibration identifiée et d'un environnement de vérification cohérent. Cette maturité relative justifie qu'il constitue le cœur technique du présent rapport.

\section{Place de la mesure temporelle dans SPADMIC}

Une matrice de \textbf{SPAD} produit des impulsions brèves, asynchrones et potentiellement nombreuses. Pour rendre ces impulsions exploitables, l'électronique de lecture doit remplir plusieurs rôles simultanément :
\begin{itemize}
    \item détecter l'événement et le convertir en information numérique ;
    \item l'associer à un instant de référence avec une finesse compatible avec l'application visée ;
    \item conserver une plage dynamique suffisante pour couvrir des délais utiles allant bien au-delà du seul pas de quantification ;
    \item transmettre les données sans imposer un volume de sortie ou un nombre de broches irréaliste.
\end{itemize}

Cette combinaison de contraintes explique pourquoi un \textbf{TDC} destiné à \textbf{SPADMIC} ne peut pas être conçu comme un bloc isolé, uniquement optimisé pour la résolution. Il doit être pensé comme une brique mixte, à l'interface entre un front-end asynchrone, une logique de capture rapide, des mécanismes de stockage temporaires et une chaîne de lecture numérique cadencée par une horloge système bien plus lente.

\begin{figure}[H]
    \centering
    \includegraphics[width=0.98\textwidth]{tdc/mptdc_overview.pdf}
    \caption{Vue fonctionnelle simplifiée du bloc \textbf{MPTDC} étudié dans ce rapport. Cette représentation restera volontairement globale à ce stade ; la hiérarchie détaillée sera reprise au chapitre \ref{chap:architecture-mptdc}.}
    \label{fig:mptdc-overview-context}
\end{figure}

La figure \ref{fig:mptdc-overview-context} montre bien cette dualité. Le cœur temporel reste fondé sur le principe Vernier, mais il est immédiatement entouré de fonctions de contrôle, de capture et de lecture. Dans le cadre de \textbf{SPADMIC}, la valeur du \textbf{TDC} provient donc autant de sa résolution temporelle que de sa capacité à s'insérer proprement dans une chaîne de traitement plus large.

\section{Contraintes fonctionnelles et technologiques}

\subsection{Contraintes liées à la mesure}

Le premier niveau de contrainte est directement lié à la mesure elle-même. Un \textbf{TDC} intégré dans \textbf{SPADMIC} doit combiner :
\begin{itemize}
    \item une résolution suffisamment fine pour rester pertinente dans un contexte de photon timing ;
    \item une plage dynamique utile, bien supérieure à quelques pas élémentaires ;
    \item une latence de réarmement courte, afin de ne pas transformer la chaîne de lecture en goulet d'étranglement ;
    \item une capacité à gérer plusieurs événements ou plusieurs détections associées à une même conversion ;
    \item une forme de robustesse face aux variations de procédé, de tension et de température.
\end{itemize}

Ces exigences ne peuvent pas être satisfaites par une architecture purement idéale. Dès que l'on passe à une implémentation réelle, il faut tenir compte du jitter des oscillateurs, de la non-uniformité de la grille fine, des frontières de comptage, des contraintes de communication entre domaines d'horloge et des limites imposées par la lecture des données.

\subsection{Contraintes d'intégration dans un circuit en \textit{ASIC}}

Le second niveau de contrainte provient de l'intégration elle-même. Le projet vise une implémentation en technologie \textbf{XFAB XH018} (\SI{180}{\nano\meter} CMOS), ce qui impose de conserver un regard réaliste sur la complexité numérique embarquée. Dans un tel contexte, chaque choix architectural a un impact direct sur :
\begin{itemize}
    \item la surface occupée par la matrice de détection de phase ;
    \item la profondeur des pipelines nécessaires à la fermeture temporelle ;
    \item la consommation associée aux oscillateurs, aux compteurs et à la lecture ;
    \item la difficulté de synthèse, de contraintes STA et de revue CDC ;
    \item le coût en broches et en bande passante de la sortie numérique.
\end{itemize}

Autrement dit, une bonne architecture n'est pas seulement celle qui donne la meilleure formule de conversion. C'est aussi celle qui reste implémentable, vérifiable et intégrable sans déplacer les difficultés vers un autre niveau du système.

\section{Pourquoi le rapport se concentre sur le TDC}

L'architecture globale de \textbf{SPADMIC} n'étant pas encore entièrement figée, il serait artificiel de présenter l'ensemble du chip comme un objet parfaitement stabilisé. À l'inverse, le bloc de conversion temps-numérique dispose déjà :
\begin{itemize}
    \item d'un socle théorique solidement documenté dans la littérature et dans la thèse de référence \cite{annagrebah2021thesis} ;
    \item d'une implémentation active suffisamment riche pour révéler de vraies contraintes d'ingénierie ;
    \item d'une infrastructure de simulation, de collecte et de calibration permettant d'aller au-delà d'une simple description conceptuelle.
\end{itemize}

Le choix retenu dans ce rapport consiste donc à suivre un chemin volontairement descendant :
\begin{enumerate}
    \item rappeler pourquoi l'approche Vernier multiphase est pertinente sur le plan théorique ;
    \item montrer comment cette approche doit être adaptée lorsqu'elle devient un bloc RTL réel ;
    \item détailler les compromis de mise en œuvre qui font passer d'une architecture de référence à une brique intégrable dans \textbf{SPADMIC}.
\end{enumerate}

Cette méthode évite deux écueils. Le premier serait de rester à un niveau trop académique, sans montrer les difficultés réelles de l'implémentation. Le second serait d'entrer directement dans le RTL sans avoir explicité les raisons profondes des choix retenus.

\section{Objectifs techniques retenus pour le bloc étudié}

Dans l'état actuel du projet, les objectifs poursuivis par le bloc \textbf{MPTDC} peuvent être résumés de la manière suivante :
\begin{itemize}
    \item conserver le principe Vernier à deux oscillateurs, de façon à bénéficier d'une résolution issue d'une différence de délais plutôt que d'un délai absolu ;
    \item exploiter une observation multiphase sur une matrice \mbox{$9\times9$}, afin d'obtenir une granularité fine tout en gardant une surface maîtrisable ;
    \item accepter des impulsions \texttt{START}/\texttt{STOP} réellement asynchrones, ce qui correspond à la nature même des événements à mesurer ;
    \item séparer la capture rapide de la lecture système au moyen de deux contextes, afin de réduire l'impact direct du drain sur le réarmement ;
    \item exporter des observables riches, pour permettre une calibration \textit{offline} plus souple et plus évolutive qu'une correction figée dans le silicium ;
    \item limiter le coût d'intégration par une interface de sortie étroite (\SI{16}{\bit}) et un protocole de paquetisation explicite.
\end{itemize}

Ces objectifs montrent déjà que la question centrale n'est pas seulement : \og comment obtenir quelques picosecondes de résolution ? \fg{} La vraie question est plus large : \og comment construire un \textbf{TDC} de haute précision qui reste cohérent avec une implémentation \textit{ASIC}, une stratégie de vérification et un budget d'intégration réaliste ? \fg{}

\begin{summaryBox}
Le bloc étudié dans ce rapport n'est pas présenté comme un démonstrateur théorique isolé. Il est analysé comme une brique d'architecture réelle, soumise à des contraintes de technologie, de domaines d'horloge, de lecture, de surface et de calibration. C'est cette articulation entre théorie et implémentation qui justifie le niveau de détail retenu dans les chapitres suivants.
\end{summaryBox}

\clearpage
